\section{Strengths and Weaknesses}\label{sec:sw}

\subsection{Strengths}

\begin{itemize}
\item \textbf{One geometry engine, four answers.} All four questions consume the same
closed-set predicate layer (clearance, tangent-blocked line of sight, elevation cone)
and the same obstacle-aware router, so conventions are consistent and one 51-check
regression suite certifies them.
\item \textbf{Closed-form cores with certified numerics.} Q1's feasible arcs and Q2(1)'s
RMSE law are exact; grids and Monte Carlo only \emph{verify}, and every random
experiment is seed-fixed and reproducible.
\item \textbf{Bounds reported honestly.} Q2(1) attains its Fisher lower bound exactly;
Q2(2)'s optimality gap ($7.17\,\mathrm{s}$; $6.78\,\mathrm{s}$ against the refined
incumbent) is stated rather than disguised, and the restricted structure space is
declared as a convention.
\item \textbf{Bounded, replayable schedules.} Q2(2) and Q4 results are simulator
replays, not paper arithmetic; route legality is certified by an independent
segment-level auditor with explicit gap numbers.
\item \textbf{Independently cross-validated.} A second, independently written team
implementation reproduces every Q1--Q3 headline exactly (same Q1 station, Q2(1) pair
and RMSE, Q3 $s_{\max}$) and refines the Q2(2) incumbent to $10.651\,\mathrm{s}$ on the
same structure family---a replay-verified $3.6\%$ bound on station-level
suboptimality; Q4's value boundary is unchanged.
\item \textbf{Robustness as budgets, not anecdotes.} Sensitivity is expressed as full
feasible regions (Q3's $(a,b)$ map) or swept parameters ($t_{rec}$, $\pm10\%$
kinematics), each with a quantified margin.
\end{itemize}

\subsection{Weaknesses}

\begin{itemize}
\item \textbf{Restricted structure space (Q2(2)/Q4).} We search single-doubler
structures with at most two task slots per vehicle (360 structures); a multi-doubler
plan could in principle close part of the $7.17\,\mathrm{s}$ gap, though no explored
structure beats the incumbent and the station-level residual is bounded at $3.6\%$
by the cross-implementation search.
\item \textbf{Visibility-graph conservatism.} Corner-vertex paths can exceed the true
geodesic optimum by up to $0.643\,\mathrm{m}$ per corner, so all travel times are
feasible upper bounds (the independent engine's exact $R_V$-neighbourhood predicate
agrees on every audited leg).
\item \textbf{Linearized pointing error (Q3).} The small-angle chain is validated to
$0.19\%$ relative error; at $s_{\max}$ the exact $\arctan$ law gives $0.95040$, so the
linearized cap is \emph{conservative} by $0.17\,\mathrm{m}$. A field safety margin is
still recommended, since hand-recomputable budgets should not sit exactly on the
boundary.
\item \textbf{Independence assumptions.} The sensor errors are modeled as mutually
independent Gaussians, as the problem states; correlated or heavy-tailed sensors would
require the distribution-free route, which collapses the travel budget
($2.45\,\mathrm{m}$ under Chebyshev).
\item \textbf{Greedy P2 is a fallback, not a guarantee.} The rule-based dispatcher
stays within $4.53\,\mathrm{s}$ of P1 on the tested families but carries no optimality
certificate beyond that empirical range.
\end{itemize}
